% Options for packages loaded elsewhere
\PassOptionsToPackage{unicode}{hyperref}
\PassOptionsToPackage{hyphens}{url}
%
\documentclass[
]{article}
\usepackage{amsmath,amssymb}
\usepackage{iftex}
\ifPDFTeX
  \usepackage[T1]{fontenc}
  \usepackage[utf8]{inputenc}
  \usepackage{textcomp} % provide euro and other symbols
\else % if luatex or xetex
  \usepackage{unicode-math} % this also loads fontspec
  \defaultfontfeatures{Scale=MatchLowercase}
  \defaultfontfeatures[\rmfamily]{Ligatures=TeX,Scale=1}
\fi
\usepackage{lmodern}
\ifPDFTeX\else
  % xetex/luatex font selection
\fi
% Use upquote if available, for straight quotes in verbatim environments
\IfFileExists{upquote.sty}{\usepackage{upquote}}{}
\IfFileExists{microtype.sty}{% use microtype if available
  \usepackage[]{microtype}
  \UseMicrotypeSet[protrusion]{basicmath} % disable protrusion for tt fonts
}{}
\makeatletter
\@ifundefined{KOMAClassName}{% if non-KOMA class
  \IfFileExists{parskip.sty}{%
    \usepackage{parskip}
  }{% else
    \setlength{\parindent}{0pt}
    \setlength{\parskip}{6pt plus 2pt minus 1pt}}
}{% if KOMA class
  \KOMAoptions{parskip=half}}
\makeatother
\usepackage{xcolor}
\usepackage[margin=1in]{geometry}
\usepackage{color}
\usepackage{fancyvrb}
\newcommand{\VerbBar}{|}
\newcommand{\VERB}{\Verb[commandchars=\\\{\}]}
\DefineVerbatimEnvironment{Highlighting}{Verbatim}{commandchars=\\\{\}}
% Add ',fontsize=\small' for more characters per line
\usepackage{framed}
\definecolor{shadecolor}{RGB}{248,248,248}
\newenvironment{Shaded}{\begin{snugshade}}{\end{snugshade}}
\newcommand{\AlertTok}[1]{\textcolor[rgb]{0.94,0.16,0.16}{#1}}
\newcommand{\AnnotationTok}[1]{\textcolor[rgb]{0.56,0.35,0.01}{\textbf{\textit{#1}}}}
\newcommand{\AttributeTok}[1]{\textcolor[rgb]{0.13,0.29,0.53}{#1}}
\newcommand{\BaseNTok}[1]{\textcolor[rgb]{0.00,0.00,0.81}{#1}}
\newcommand{\BuiltInTok}[1]{#1}
\newcommand{\CharTok}[1]{\textcolor[rgb]{0.31,0.60,0.02}{#1}}
\newcommand{\CommentTok}[1]{\textcolor[rgb]{0.56,0.35,0.01}{\textit{#1}}}
\newcommand{\CommentVarTok}[1]{\textcolor[rgb]{0.56,0.35,0.01}{\textbf{\textit{#1}}}}
\newcommand{\ConstantTok}[1]{\textcolor[rgb]{0.56,0.35,0.01}{#1}}
\newcommand{\ControlFlowTok}[1]{\textcolor[rgb]{0.13,0.29,0.53}{\textbf{#1}}}
\newcommand{\DataTypeTok}[1]{\textcolor[rgb]{0.13,0.29,0.53}{#1}}
\newcommand{\DecValTok}[1]{\textcolor[rgb]{0.00,0.00,0.81}{#1}}
\newcommand{\DocumentationTok}[1]{\textcolor[rgb]{0.56,0.35,0.01}{\textbf{\textit{#1}}}}
\newcommand{\ErrorTok}[1]{\textcolor[rgb]{0.64,0.00,0.00}{\textbf{#1}}}
\newcommand{\ExtensionTok}[1]{#1}
\newcommand{\FloatTok}[1]{\textcolor[rgb]{0.00,0.00,0.81}{#1}}
\newcommand{\FunctionTok}[1]{\textcolor[rgb]{0.13,0.29,0.53}{\textbf{#1}}}
\newcommand{\ImportTok}[1]{#1}
\newcommand{\InformationTok}[1]{\textcolor[rgb]{0.56,0.35,0.01}{\textbf{\textit{#1}}}}
\newcommand{\KeywordTok}[1]{\textcolor[rgb]{0.13,0.29,0.53}{\textbf{#1}}}
\newcommand{\NormalTok}[1]{#1}
\newcommand{\OperatorTok}[1]{\textcolor[rgb]{0.81,0.36,0.00}{\textbf{#1}}}
\newcommand{\OtherTok}[1]{\textcolor[rgb]{0.56,0.35,0.01}{#1}}
\newcommand{\PreprocessorTok}[1]{\textcolor[rgb]{0.56,0.35,0.01}{\textit{#1}}}
\newcommand{\RegionMarkerTok}[1]{#1}
\newcommand{\SpecialCharTok}[1]{\textcolor[rgb]{0.81,0.36,0.00}{\textbf{#1}}}
\newcommand{\SpecialStringTok}[1]{\textcolor[rgb]{0.31,0.60,0.02}{#1}}
\newcommand{\StringTok}[1]{\textcolor[rgb]{0.31,0.60,0.02}{#1}}
\newcommand{\VariableTok}[1]{\textcolor[rgb]{0.00,0.00,0.00}{#1}}
\newcommand{\VerbatimStringTok}[1]{\textcolor[rgb]{0.31,0.60,0.02}{#1}}
\newcommand{\WarningTok}[1]{\textcolor[rgb]{0.56,0.35,0.01}{\textbf{\textit{#1}}}}
\usepackage{graphicx}
\makeatletter
\def\maxwidth{\ifdim\Gin@nat@width>\linewidth\linewidth\else\Gin@nat@width\fi}
\def\maxheight{\ifdim\Gin@nat@height>\textheight\textheight\else\Gin@nat@height\fi}
\makeatother
% Scale images if necessary, so that they will not overflow the page
% margins by default, and it is still possible to overwrite the defaults
% using explicit options in \includegraphics[width, height, ...]{}
\setkeys{Gin}{width=\maxwidth,height=\maxheight,keepaspectratio}
% Set default figure placement to htbp
\makeatletter
\def\fps@figure{htbp}
\makeatother
\setlength{\emergencystretch}{3em} % prevent overfull lines
\providecommand{\tightlist}{%
  \setlength{\itemsep}{0pt}\setlength{\parskip}{0pt}}
\setcounter{secnumdepth}{-\maxdimen} % remove section numbering
\usepackage{mathpazo}
\usepackage{booktabs}
\usepackage{rotating}
\usepackage{makecell}
\usepackage{hyperref}
%% \usepackage[a4paper,margin=2cm]{geometry}


\newcommand{\orcidlogo}{\textcolor{orcidlogocol}{\aiOrcid}}
\usepackage[font=Large,labelfont=bf,textfont=bf]{caption}
\DeclareCaptionLabelFormat{addC}{#1 C#2}


%% This solution for adding the ORCID logo was taken from:
%% https://tex.stackexchange.com/questions/445563/ieeetran-how-to-include-orcid-in-tex-pdf-with-pdflatex/445583
%% and adapted simply by scaling the icon a little smaller than | to +
\usepackage{scalerel}
\usepackage{tikz}
\usetikzlibrary{svg.path}

\definecolor{orcidlogocol}{HTML}{A6CE39}
\tikzset{
  orcidlogo/.pic={
    \fill[orcidlogocol] svg{M256,128c0,70.7-57.3,128-128,128C57.3,256,0,198.7,0,128C0,57.3,57.3,0,128,0C198.7,0,256,57.3,256,128z};
    \fill[white] svg{M86.3,186.2H70.9V79.1h15.4v48.4V186.2z}
                 svg{M108.9,79.1h41.6c39.6,0,57,28.3,57,53.6c0,27.5-21.5,53.6-56.8,53.6h-41.8V79.1z M124.3,172.4h24.5c34.9,0,42.9-26.5,42.9-39.7c0-21.5-13.7-39.7-43.7-39.7h-23.7V172.4z}
                 svg{M88.7,56.8c0,5.5-4.5,10.1-10.1,10.1c-5.6,0-10.1-4.6-10.1-10.1c0-5.6,4.5-10.1,10.1-10.1C84.2,46.7,88.7,51.3,88.7,56.8z};
  }
}

\newcommand\orcidicon[1]{\href{https://orcid.org/#1}{\mbox{\scalerel*{
\begin{tikzpicture}[yscale=-1,transform shape]
\pic{orcidlogo};
\end{tikzpicture}
}{+}}}}

%% Set hyperlinks in colour to the "CODECHECK colour".
\hypersetup{colorlinks=true, urlcolor = [rgb]{0,.5,0.2}}
\ifLuaTeX
  \usepackage{selnolig}  % disable illegal ligatures
\fi
\IfFileExists{bookmark.sty}{\usepackage{bookmark}}{\usepackage{hyperref}}
\IfFileExists{xurl.sty}{\usepackage{xurl}}{} % add URL line breaks if available
\urlstyle{same}
\hypersetup{
  pdftitle={CODECHECK certificate 2023-011},
  hidelinks,
  pdfcreator={LaTeX via pandoc}}

\title{CODECHECK certificate 2023-011}
\usepackage{etoolbox}
\makeatletter
\providecommand{\subtitle}[1]{% add subtitle to \maketitle
  \apptocmd{\@title}{\par {\large #1 \par}}{}{}
}
\makeatother
\subtitle{\url{https://doi.org/10.5281/zenodo.FIXME}}
\author{}
\date{\vspace{-2.5em}}

\begin{document}
\maketitle

\centerline{\includegraphics[width=4cm]{/Users/stephen/NOBACKUP/RLIB/codecheck/extdata/codecheck_logo.pdf}}\vspace*{2cm}

\begin{table}[ht]
\centering
\begin{tabular}{lp{10cm}}
  \hline
Item & Value \\ 
  \hline
Title & Does enforcing glenohumeral joint stability matter? A new rapid muscle redundancy solver highlights the importance of non-superficial shoulder muscles
 \\ 
  Authors & Italo Belli \orcidicon{0000-0002-8064-2251} , Ajay Seth \orcidicon{0000-0003-4217-1580}  \\ 
  Reference & BioRxiv preprint (2003)  \url{https://www.biorxiv.org/content/10.1101/2023.07.11.548542v1}
 \\ 
  Codechecker & Stephen J. Eglen \orcidicon{0000-0001-8607-8025}  \\ 
  Date of check & 2023-09-18  13:00:00 \\ 
  Summary & Codecheck performed interactively as part of the Delft 2023 workshop.
 \\ 
  Repository & \url{https://github.com/codecheckers/Reproduction-HancockXXX} \\ 
   \hline
\end{tabular}
\caption{CODECHECK summary} 
\end{table}

\begin{table}[ht]
\centering
\begin{tabular}{p{6cm}p{6cm}p{2cm}}
  \hline
Output & Comment & Size (b) \\ 
  \hline
\href{https://github.com/codecheckers/Reproduction-HancockXXX/blob/master/codecheck/outputs/figure3-screenshot.png}{\path{codecheck/figure3-screenshot.png}} & manuscript Figure 3 (composite) & 1198909 \\ 
   \hline
\end{tabular}
\caption{Summary of output files generated} 
\end{table}

\section{Summary}\label{summary}

Source of project
\url{https://github.com/ComputationalBiomechanicsLab/rmr-solver}

This code was fairly straightforward to codecheck. {[}\ldots{} ADD MORE
INTERESTING FINDINGS HERE..{]}

\clearpage

\section{CODECHECKER notes}\label{codechecker-notes}

The GitHub repo \ldots{} This check is based on the commit
\texttt{fefcb73d0799aa32c8f9fe80d33104ced75d7fff}. Code was written in
\ldots{} I went through the following steps \ldots{} One hard problem
was \ldots{} I added the following files \ldots{} using tools/methods
\ldots{}

This took \ldots{} minutes to complete on \{a large workstation, my
laptop\}.

\subsection{Recommendations}\label{recommendations}

I suggest to the authors to consider the following suggestions for their
next publication or workflow:

\begin{itemize}
\tightlist
\item
  \ldots{}
\end{itemize}

(document here if any of the suggestions were taken up by the authors in
the meantime - do not remove any, keep track of contributions via
feedback)

\clearpage

\section{Manifest files}\label{manifest-files}

\subsection{figure3-screenshot.png}\label{figure3-screenshot.png}

\textbf{Comment:} manuscript Figure 3 (composite)

\includegraphics{/Users/stephen/proj/2020/codecheckers/rmr-solver/codecheck/outputs/figure3-screenshot.png}
\clearpage 

\clearpage

\subsection{Acknowledgements}\label{acknowledgements}

I would like to thank Dr Bhatt and his team for promptly answering any
queries I had with this reproduction. CODECHECK is financially supported
by the Mozilla foundation.

\section{Citing this document}\label{citing-this-document}

Stephen J. Eglen (2023). CODECHECK Certificate 2023-011. Zenodo.
\url{https://doi.org/10.5281/zenodo.FIXME}

\section{About CODECHECK}\label{about-codecheck}

This certificate confirms that the codechecker could independently
reproduce the results of a computational analysis given the data and
code from a third party. A CODECHECK does not check whether the original
computation analysis is correct. However, as all materials required for
the reproduction are freely available by following the links in this
document, the reader can then study for themselves the code and data.

\section{About this document}\label{about-this-document}

This document was created using \href{https://rmarkdown.rstudio.com/}{R
Markdown} using the
\href{https://github.com/codecheckers/codecheck}{\texttt{codecheck}} R
package. \texttt{make\ codecheck.pdf} will regenerate the report file.

\begin{Shaded}
\begin{Highlighting}[]
\FunctionTok{sessionInfo}\NormalTok{()}
\end{Highlighting}
\end{Shaded}

\begin{verbatim}
## R version 4.3.1 (2023-06-16)
## Platform: aarch64-apple-darwin22.4.0 (64-bit)
## Running under: macOS Ventura 13.5.2
## 
## Matrix products: default
## BLAS:   /opt/homebrew/Cellar/openblas/0.3.23/lib/libopenblasp-r0.3.23.dylib 
## LAPACK: /opt/homebrew/Cellar/r/4.3.1/lib/R/lib/libRlapack.dylib;  LAPACK version 3.11.0
## 
## locale:
## [1] en_GB.UTF-8/en_GB.UTF-8/en_GB.UTF-8/C/en_GB.UTF-8/en_GB.UTF-8
## 
## time zone: Europe/London
## tzcode source: internal
## 
## attached base packages:
## [1] stats     graphics  grDevices utils     datasets 
## [6] methods   base     
## 
## other attached packages:
##  [1] readr_2.1.3          tibble_3.1.8        
##  [3] xtable_1.8-4         yaml_2.3.6          
##  [5] rprojroot_2.0.3      knitr_1.41          
##  [7] codecheck_0.1.0.9005 git2r_0.32.0        
##  [9] parsedate_1.3.1      R.cache_0.16.0      
## [11] gh_1.3.1            
## 
## loaded via a namespace (and not attached):
##  [1] utf8_1.2.2        generics_0.1.3    xml2_1.3.3       
##  [4] stringi_1.7.12    zen4R_0.8         httpcode_0.3.0   
##  [7] hms_1.1.2         digest_0.6.31     magrittr_2.0.3   
## [10] evaluate_0.19     fastmap_1.1.0     R.oo_1.25.0      
## [13] jsonlite_1.8.4    R.utils_2.12.2    whisker_0.4.1    
## [16] DBI_1.1.3         crul_1.4.0        httr_1.4.4       
## [19] purrr_1.0.0       fansi_1.0.3       cli_3.6.0        
## [22] rlang_1.0.6       R.methodsS3_1.8.2 ellipsis_0.3.2   
## [25] cachem_1.0.6      tools_4.3.1       tzdb_0.3.0       
## [28] memoise_2.0.1     dplyr_1.0.10      curl_4.3.3       
## [31] assertthat_0.2.1  vctrs_0.5.1       R6_2.5.1         
## [34] lifecycle_1.0.3   stringr_1.5.0     fs_1.5.2         
## [37] pkgconfig_2.0.3   rorcid_0.7.0      osfr_0.2.9       
## [40] pillar_1.8.1      glue_1.6.2        xfun_0.36        
## [43] tidyselect_1.2.0  keyring_1.3.1     htmltools_0.5.4  
## [46] rmarkdown_2.19    compiler_4.3.1    fauxpas_0.5.2
\end{verbatim}

\end{document}
